%% ============================================================================
%%  Packet-Semantic Constraints for Problem-Space Adversarial Attacks
%%  Companion note to the IEEE Access resubmission (v4) of Stochastic ODE-Guard.
%%  Compile with pdflatex (self-contained).
%% ============================================================================
\documentclass[11pt,letterpaper]{article}

\usepackage[margin=1in]{geometry}
\usepackage[T1]{fontenc}
\usepackage{lmodern}
\usepackage{amsmath,amssymb}
\usepackage{booktabs}
\usepackage{array}
\usepackage{tabularx}
\usepackage{enumitem}
\usepackage{xcolor}
\usepackage[colorlinks=true,linkcolor=blue!60!black,urlcolor=blue!60!black]{hyperref}
\usepackage{fancyhdr}
\usepackage{titlesec}

\setlength{\parindent}{0pt}
\setlength{\parskip}{6pt}
\pagestyle{fancy}
\fancyhf{}
\lhead{\small SODE-Guard — v4 companion note}
\rhead{\small Packet Semantics}
\cfoot{\thepage}
\titleformat{\section}{\normalfont\Large\bfseries\color{blue!60!black}}{\thesection}{0.7em}{}

\title{\vspace{-1.2cm}\textbf{Packet-Semantic Constraints for
       Problem-Space Adversarial Attacks}\\[4pt]
       \normalsize Companion note to the SODE-Guard v4 manuscript
       (IEEE Access MS \texttt{Access-2026-27603})}
\author{R.\ N.\ Anaedevha, A.\ G.\ Trofimov, and Y.\ V.\ Borodachev\\
        \normalsize National Research Nuclear University MEPhI, Moscow, Russia}
\date{\normalsize \today}

\begin{document}
\maketitle
\thispagestyle{fancy}
\vspace{-0.5cm}

\section*{Purpose}
Reviewer~2 (concern R2.11) asked whether adversarial perturbations that
live in the 83-dimensional flow feature space can actually be realised
by network traffic. The short answer is: not all of them. Certain
features are \emph{constrained} by the packet grammar; others are
\emph{derived} from raw counts and must remain consistent with them.

This note lists the constraint families implemented by
\texttt{src/attacks/semantic/feasibility.py} in the reproducibility
repository. All constraints are applied at every step of the PGD-40 /
Carlini--Wagner adversary so that the ``problem-space'' robust
accuracy in Table~9 of the v4 manuscript uses only \emph{feasible}
perturbations.

%% ============================================================================
\section{Constraint families}
%% ============================================================================
\begin{enumerate}[leftmargin=1.4em]
  \item \textbf{Box constraints.} Every feature has a natural
        non-negative lower bound (byte counts, packet counts,
        durations, TLS extension counts) or a flag-bit domain
        $\{0,1\}$.
  \item \textbf{Integer constraints.} Packet counts, byte totals, and
        TLS enum indices (\texttt{cipher\_suite\_id},
        \texttt{pqc\_kx\_id}, \texttt{alpn\_id}, \ldots) round to the
        nearest non-negative integer.
  \item \textbf{Flag constraints.} All \texttt{flag\_*} features
        (\texttt{flag\_psh}, \texttt{flag\_syn}, \texttt{flag\_ack},
        \texttt{flag\_fin}, \ldots) are Bernoulli --- clipped and
        rounded.
  \item \textbf{Derived-ratio recomputation.} Whenever both numerator
        and denominator of a derived ratio
        (\texttt{fwd\_bwd\_byte\_ratio}, \texttt{fwd\_bwd\_pkt\_ratio},
        \texttt{down\_up\_ratio}) are present, the ratio is recomputed
        from the projected raw counts so the adversary cannot
        desynchronise it.
  \item \textbf{Protocol-structural constraints.}
        $\texttt{min\_seg\_size\_fwd}\le \texttt{avg\_fwd\_seg\_size}\le \texttt{fwd\_pkt\_len\_max}$;
        $\texttt{total\_pkts}=\texttt{fwd\_packets}+\texttt{bwd\_packets}$;
        $\texttt{handshake\_rtt}\ge 0$;
        TLS / PQC identifiers $\ge 0$.
\end{enumerate}

%% ============================================================================
\section{What the projector does not enforce}
%% ============================================================================
The projector is \emph{conservative}: it is permissive of features
that cannot be tied to a single raw value by the CICFlowMeter grammar,
such as \texttt{entropy\_payload} or the JA4 hash bucket. These are
allowed to move freely inside their box because a real attacker can,
in principle, manipulate them via payload crafting. Table~9 of the
revised manuscript reports both the unconstrained
(``feature-space'') and the constrained (``problem-space'') attack
results so the reader can see the delta.

%% ============================================================================
\section{Empirical impact}
%% ============================================================================
Under PGD-40 at $\varepsilon=0.03$ on the ICS3D test partition:

\begin{center}\small
\begin{tabular}{@{}lccc@{}}
\toprule
\textbf{Threat model} & \textbf{SDE-TGNN F1} & \textbf{SODE-Guard F1} & \textbf{Gap} \\
\midrule
Feature-space  & 0.858 & 0.922 & $+6.4$ \\
Problem-space  & 0.842 & 0.914 & $+7.2$ \\
\bottomrule
\end{tabular}
\end{center}

SODE-Guard's advantage \emph{widens} under the tighter constraints,
which is consistent with the anti-concentration argument: the
effective chaos degree of the smoothed margin is unchanged by
projecting to a smaller feasibility set, so the certificate's
$\varepsilon^{1/d^\star}$ tail decay still applies. See Proposition~5
and Corollary~1 of the v4 manuscript for the analytical statement.

%% ============================================================================
\section*{Availability}
%% ============================================================================
Implementation:
\url{https://github.com/rogerpanel/CV/tree/main/SODE-Guard/src/attacks/semantic}.

Unit tests: \texttt{tests/test\_feasibility\_projector.py} exercises
each constraint family on synthetic inputs and verifies that the
projector is idempotent.

\end{document}
